\chapter{Общий parent Mathai--Quillen Thom-мультиплета}
\label{ch:v10-h15-r2-mathai-quillen-thom-multiplet-parent}

Предыдущий гейт оставил один вопрос: может ли условная auxiliary-копия быть
получена как полный Thom-мультиплет, а не как одиночное Gaussian-поле?
Начнём с положительного операторного результата. Для relative projector
$P_{\rm rel}$ suspension swap не является новым выбором:
\begin{equation}
 S_{\rm sus}=I_{16}-2P_{\rm rel}.
\end{equation}
Если $P_\pm=(I_{16}\pm\Gamma_{\rm sus})/2$, то его ориентированная часть
равна
\begin{equation}
 d_{\rm car}=P_-S_{\rm sus}P_+
 =\begin{pmatrix}0&0\\I_8&0\end{pmatrix},
 \qquad d_{\rm car}^2=0,
 \qquad \rank d_{\rm car}=8.
\end{equation}
Следовательно, carrier differential действительно выводится из уже
построенных Hodge- и grading-данных, условно на их admission.

Однако полный Mathai--Quillen объект содержит квартет
\begin{equation}
 (\Sigma,\psi_\Sigma,\chi_\Sigma,H_\Sigma),
\end{equation}
где $\Sigma,H_\Sigma$ бозонны, а $\psi_\Sigma,\chi_\Sigma$ имеют нечётную
статистику. На 32-мерном вещественном пространстве его differential имеет
два contractible блока,
\begin{equation}
 \delta\Sigma=\psi_\Sigma,
 \quad \delta\psi_\Sigma=0,
 \quad \delta\chi_\Sigma=H_\Sigma,
 \quad \delta H_\Sigma=0.
\end{equation}
Точная матрица $\delta$ имеет rank $16$, квадрат ноль и
$\ker\delta=\operatorname{im}\delta$. Она нечётна по field parity и
коммутирует с gauge- и Real-действиями.

Полная мера также условно замыкается. После стандартного сдвига auxiliary
contour бозонный квадратичный оператор равен
\begin{equation}
 \mathcal H_B=Q^2\oplus I_8,
\end{equation}
а нечётный билинейный оператор --- $Q$ в паре
$\chi_\Sigma Q\psi_\Sigma$. Поэтому
\begin{equation}
 \det\mathcal H_B=(\det Q)^2=3969^2,
 \qquad \det Q=3969,
\end{equation}
и нормированный Gaussian--Berezin determinant ratio точно равен единице.
Таким образом, полный условный квартет не создаёт дополнительного
детерминантного масштаба и не оставляет on-shell cohomology.

Строгий physical ledger, однако, отрицателен. Текущий parent содержит
только физическую $\Sigma$, rank $8$. Предыдущая условная достройка даёт
две бозонные координаты rank $16$, тогда как полный квартет требует
$32$ вещественных направлений. Inherited injection двух нечётных полей
имеет rank ноль; дефицит равен шестнадцати направлениям. Градуировка
carrier не превращает коэффициенты поля в Grassmann-переменные и не
порождает Berezin measure.

\begin{theorem}[Условное замыкание полного Thom-квартета]
Relative-Hodge projector и suspension grading канонически выводят
nilpotent carrier differential. Полный Mathai--Quillen-квартет условно
имеет нулевую cohomology, положительную бозонную меру и точное сокращение
Gaussian--Berezin determinants. Conditional status равен $7/8$.
\end{theorem}

\begin{nogocriterion}[Carrier parity не является статистикой поля]
Нельзя считать нечётные Thom-поля унаследованными только потому, что
конечномерный carrier имеет grading. Без независимых Grassmann-координат,
nilpotent derivation на алгебре полей и Berezin measure физический статус
остаётся $3/8$; недостающий odd rank равен $16$.
\end{nogocriterion}

\begin{protocol}\begin{enumerate}
\item derived oriented carrier edge rank: \textbf{$8$};
\item full Thom differential rank/cohomology: \textbf{$16/0$};
\item required full field dimension: \textbf{$32$};
\item inherited/conditional bosonic rank: \textbf{$8/16$};
\item inherited odd rank and deficit: \textbf{$0/16$};
\item bosonic Hessian inertia: \textbf{$(0_-,0_0,16_+)$};
\item fermionic determinant: \textbf{$3969$};
\item normalized determinant ratio: \textbf{$1$};
\item conditional/inherited status: \textbf{$7/8$ и $3/8$};
\item следующий гейт: \textbf{аудит происхождения статистики odd Thom-пары}.
\end{enumerate}\end{protocol}

Машинный сертификат сохранён в
\path{s2t_v10_particle_wrinkle_dislocation_callias_equal_charge_twist_m4_h15_r2_mathai_quillen_thom_multiplet_common_parent_origin_gate_results.json}.